\documentclass[12pt,a4paper]{article}
\usepackage[utf8]{inputenc}
\usepackage[T1]{fontenc}
\usepackage{lmodern}
\usepackage{geometry}
\geometry{margin=2.5cm}
\usepackage{setspace}
\onehalfspacing
\usepackage{natbib}
\usepackage{hyperref}
\hypersetup{colorlinks=true,linkcolor=blue,citecolor=blue,urlcolor=blue}
\usepackage{csquotes}
\usepackage{amsmath,amssymb}

\title{The Floating Signifier of AI:\\Legitimacy, Authenticity, and the Moral Economy of Writing in the Generative Age}
\author{Anonymous}
\date{\today}

\begin{document}
\maketitle

\begin{abstract}
This paper examines the sociological and philosophical dimensions of the contemporary practice of labeling disliked or stylistically unfamiliar texts as ``AI-generated'' or ``AI-like.'' Drawing on classical theories of legitimacy (Weber), communicative action and validity claims (Habermas), and the concept of the empty or floating signifier (Laclau), it argues that ``AI'' has become a flexible moral and epistemic category used to contest the legitimacy of texts independently of actual generative processes. The analysis is grounded in recent observations from Japanese online spaces (notably Hatena Anonymous Diary) and broader academic and public discourse. It further addresses the paradox of hypocrisy in AI attitudes and the implications for the moral economy of writing. The paper concludes that the crisis is not primarily technological but one of legitimacy and intersubjective recognition in the age of probabilistic text generation.
\end{abstract}

\section{Introduction}

The rapid diffusion of large language models has not merely altered the production of text; it has reconfigured the conditions under which texts are granted or denied legitimacy. A striking feature of contemporary online and academic discourse is the use of ``AI'' as a pejorative label applied to writing that is perceived as overly structured, impersonal, lengthy, or otherwise aesthetically or politically unpalatable, often without evidence of actual machine generation. This practice is especially visible in spaces such as Hatena Anonymous Diary (増田), where readers routinely dismiss posts with the remark that they ``sound like AI,'' even when the stylistic features in question are consistent with trained professional or academic registers.

This paper takes as its central object the concept of \emph{legitimacy} (Weber's \emph{Legitimität}, often rendered in Japanese as 正当性 or レジティマシー). It asks: under what conditions is a text recognized as a legitimate contribution to discourse, and how has the introduction of generative AI transformed the criteria of that recognition? The argument proceeds in four stages. First, it reconstructs the classical sociological and philosophical resources necessary for analyzing legitimacy in communicative settings. Second, it shows how ``AI'' functions as a floating signifier in the sense developed by Laclau, allowing diverse grievances to be condensed into a single, morally charged category. Third, it examines the empirical contours of this phenomenon, with particular attention to the Japanese case and to documented instances of false accusation. Fourth, it addresses the structural hypocrisy that characterizes much public discussion of AI use. The paper closes by considering the implications for the future of textual authority.

\section{Theoretical Framework: Legitimacy, Validity Claims, and Empty Signifiers}

\subsection{Weber on Legitimacy}

Max Weber's account of legitimacy remains the foundational reference. For Weber, legitimacy is not a property of an order itself but a belief on the part of those subject to it that the order is valid. He distinguishes three pure types of legitimate domination: traditional, charismatic, and rational-legal \citep{weber1978}. In the domain of discourse and authorship, the relevant form is closest to the rational-legal type, in which authority rests on the belief in the legality of enacted rules and the competence of those who enact them. Yet even here, Weber insists that legality alone does not generate legitimacy; the procedures must themselves be experienced as rightful \citep{weber1978,habermas1984}.

Applied to writing, this implies that the recognition of a text as a legitimate contribution depends not only on its formal compliance with genre norms but on a background belief that its author has the standing to speak. When that standing is placed in doubt—by the suggestion that the author has outsourced the labor of composition to a machine—the text's legitimacy is simultaneously impugned.

\subsection{Habermas and Validity Claims}

Jürgen Habermas extends and radicalizes this insight in his theory of communicative action. For a speech act to contribute to rational discourse, it must redeem three validity claims: truth (Wahrheit), truthfulness or sincerity (Wahrhaftigkeit), and rightness or legitimacy (Richtigkeit) \citep{habermas1984,habermas1987}. The third claim is especially pertinent. Rightness refers to the normative appropriateness of the utterance within a given context of interaction. When interlocutors refuse to engage with the propositional content of a text and instead contest the legitimacy of its mode of production, they are effectively denying the rightness claim without addressing the others.

Recent work has applied Habermasian categories to AI discourse itself. Studies of existential-risk narratives and of large language models as communicative agents have shown that the validity claims of AI-generated speech are structurally impaired: the systems lack intentionality and therefore cannot fully redeem the claim to sincerity \citep{talking2024,communicative2026}. The present paper inverts the direction of analysis. It asks not whether AI can participate in communicative action, but how human actors use the category of AI to withhold recognition from other humans.

\subsection{Laclau and the Floating Signifier}

The linguistic vehicle of this withholding is what Ernesto Laclau terms the empty or floating signifier. An empty signifier is a term whose particular content has been emptied so that it can serve as the nodal point for a chain of equivalences among otherwise heterogeneous demands \citep{laclau1996,laclau2005}. ``AI'' currently performs this function. It can designate actual machine generation, stylistically formal prose, lengthy argumentation, the absence of personal anecdote, the use of half-width spaces in Japanese text, or simply any writing the reader finds uncongenial. Because the signifier is underdetermined, it can be attached to a wide range of grievances without requiring the speaker to specify the precise grounds of objection.

This mechanism is familiar from the literature on moral panics, in which floating signifiers such as ``woke'' or ``elite'' condense diverse anxieties into a single, mobilizable enemy \citep{prindle2023}. The innovation in the present case is that the floating signifier is attached not to a social group but to a mode of textual production, thereby converting aesthetic or intellectual disagreement into a question of moral and epistemic legitimacy.

\section{Empirical Observations}

\subsection{The Japanese Case: Hatena Anonymous Diary}

In late September 2026, Hatena Anonymous Diary introduced experimental connectors that allowed users to search, read, and post via ChatGPT and Claude. Within days, the volume of posts labeled ``AI'' in the comments rose sharply. A widely discussed entry of 3 October 2026 observed that the judgment ``this sounds like AI'' had become a convenient device for terminating discussion without engaging the argument \citep{anond20261003}. Features cited as diagnostic included excessive structure, frequent paragraph breaks, the presence of bullet points, length, and even brevity. The label thus functioned as an all-purpose dismissal.

This pattern is not unique to Japan. Parallel observations have been recorded on English-language platforms and in academic settings, where stylometric detectors and human readers alike have produced false positives against neurodivergent writers and non-native speakers whose natural registers overlap with the statistical signatures of large language models \citep{gradpilot2026,bloomberg2024,newby2026}.

\subsection{False Positives and the Crisis of Recognition}

Documented cases illustrate the practical consequences. Moira Olmsted, an autistic student, received a zero on an assignment whose formulaic style was flagged by a detector; Orion Newby, also autistic and enrolled in a university support program, had a Turnitin score of 100 percent overturned by a New York court that described the university's determination as ``without valid basis and devoid of reason'' \citep{newby2026,gradpilot2026}. These episodes demonstrate that the legitimacy of a text can be stripped away by an automated or informal judgment that the writing ``looks like AI,'' even when the author can produce evidence of independent composition.

University teaching centers have begun to warn faculty that neurodivergent students are especially vulnerable to such misclassification because traits such as repetitive phrasing and consistent structure are misread as machine output \citep{nebraska2024}. The sociological implication is clear: the criteria of textual legitimacy have become entangled with statistical proxies that systematically disadvantage certain styles of human expression.

\section{Hypocrisy and the Dual Attitude Structure}

A further complication is the widespread discrepancy between public condemnation of AI use and private reliance upon it. Surveys and qualitative reports document that many who criticize others for using generative tools themselves employ them for drafting, editing, or idea generation \citep{bbc2026,dartmouth2026}. This is not mere inconsistency; it is a structured dual attitude. Explicitly, AI use is framed as a violation of sincerity and effort; implicitly, it is accepted as a practical necessity.

Psychological research on moral hypocrisy and on the sycophantic tendencies of language models provides a partial explanation. Individuals readily rationalize their own use while applying stricter standards to others, particularly when the other's use threatens status or appears to reduce labor \citep{hallsworth2026}. In Habermasian terms, the rightness claim is selectively enforced: it is invoked against those whose texts one wishes to exclude, and suspended for one's own production.

\section{Conclusion}

The phenomena examined in this paper indicate that the central problem introduced by generative AI is not the technical difficulty of detection but the instability of legitimacy in textual communication. ``AI'' has become a floating signifier that permits the denial of recognition without the burden of substantive argument. Classical theories of legitimacy and communicative rationality remain indispensable for diagnosing this development, even as they must be extended to account for probabilistic text generation and the strategic use of technological categories in ordinary discourse.

Future research should track the stabilization or further diffusion of this signifier across linguistic communities and institutional settings. In the interim, the practical recommendation is modest: claims that a text is illegitimate because it ``sounds like AI'' should be treated as requiring the same justification as any other challenge to a validity claim. Absent such justification, they function less as epistemic judgments than as exercises in boundary maintenance.

\bibliographystyle{plainnat}
\begin{thebibliography}{99}

\bibitem[Weber(1978)]{weber1978}
Weber, Max. 1978. \emph{Economy and Society}. Edited by Guenther Roth and Claus Wittich. Berkeley: University of California Press.

\bibitem[Habermas(1984)]{habermas1984}
Habermas, Jürgen. 1984. \emph{The Theory of Communicative Action, Volume 1: Reason and the Rationalization of Society}. Translated by Thomas McCarthy. Boston: Beacon Press.

\bibitem[Habermas(1987)]{habermas1987}
Habermas, Jürgen. 1987. \emph{The Theory of Communicative Action, Volume 2: Lifeworld and System}. Translated by Thomas McCarthy. Boston: Beacon Press.

\bibitem[Laclau(1996)]{laclau1996}
Laclau, Ernesto. 1996. \emph{Emancipation(s)}. London: Verso.

\bibitem[Laclau(2005)]{laclau2005}
Laclau, Ernesto. 2005. \emph{On Populist Reason}. London: Verso.

\bibitem[Anond(2026)]{anond20261003}
Anonymous. 2026. ``AI増田判定、便利すぎて本文を読まなくて済むようになってる.'' \emph{はてな匿名ダイアリー}, 3 October. \url{https://anond.hatelabo.jp/20261003152005}.

\bibitem[GradPilot(2026)]{gradpilot2026}
GradPilot. 2026. ``AI Detectors Flag Autistic and ADHD Writers — The Evidence.'' 19 June.

\bibitem[Newby(2026)]{newby2026}
\textit{Matter of Newby v. Adelphi Univ.}, 2026 NY Slip Op 26021 (N.Y. Sup. Ct. Nassau County, Jan. 28, 2026).

\bibitem[Bloomberg(2024)]{bloomberg2024}
Bloomberg Businessweek. 2024. ``Do AI Detectors Work? Students Face False Cheating Accusations.'' 18 October.

\bibitem[Nebraska(2024)]{nebraska2024}
University of Nebraska–Lincoln Teaching Center. 2024. Guidance on AI Detection and Neurodivergent Students.

\bibitem[BBC(2026)]{bbc2026}
BBC. 2026. ``Are we all AI hypocrites?'' \emph{The Artificial Human}, 24 June.

\bibitem[Dartmouth(2026)]{dartmouth2026}
New York Times. 2026. ``College Leaders Are Using A.I. Students Call It Hypocrisy.'' 29 September.

\bibitem[Hallsworth(2026)]{hallsworth2026}
Hallsworth, Michael. 2026. Remarks on effort and value in AI use, cited in BBC, ``Are we all AI hypocrites?''

\bibitem[Talking(2024)]{talking2024}
``Talking existential risk into being: a Habermasian critical discourse perspective to AI hype.'' \emph{AI and Ethics}, 2024.

\bibitem[Communicative(2026)]{communicative2026}
``Communicative Action in the Age of Artificial Intelligence: Theorizing LLMs as Communicative and Strategic Agents.'' \emph{Philosophy \& Technology}, 2026.

\bibitem[Prindle(2023)]{prindle2023}
Prindle Institute for Ethics. 2023. ``Indeterminate Fear: Moral Panics and Floating Signifiers.'' 9 August.

\bibitem[Liang et al.(2023)]{liang2023}
Liang, Weixin, et al. 2023. ``GPT Detectors Are Biased Against Non-Native English Writers.'' \emph{Patterns} 4 (7).

\bibitem[Chambers and Kelley(2026)]{chambers2026}
Chambers, Summer, and Matthew C. Kelley. 2026. ``The Misclassification of Autistic Writing as AI-Generated.'' arXiv:2607.14729.

\bibitem[Westenberg(2026)]{westenberg2026}
Westenberg, JA. 2026. ``AI Writing Witch Hunts Hurt Autistic Writers Most.'' \emph{Boing Boing}, 6 April.

\bibitem[Sadasivan et al.(2023)]{sadasivan2023}
Sadasivan, Vinu Sankar, et al. 2023. ``Can AI-Generated Text be Reliably Detected?'' arXiv:2303.11156.

\bibitem[OpenAI(2023)]{openai2023}
OpenAI. 2023. ``New AI Classifier for Indicating AI-Written Text.'' (Discontinued.)

\bibitem[Turnitin(2024)]{turnitin2024}
Turnitin. 2024. Documentation on AI Writing Detection False Positive Rates.

\bibitem[Stanford HAI(2023)]{stanford2023}
Stanford HAI. 2023. ``AI Detectors Biased Against Non-Native English Writers.''

\bibitem[Nature(2026)]{nature2026}
``On the Troubling Rise of Generative AI Suspicion in Academic Publishing.'' \emph{Nature Machine Intelligence}, 2026.

\bibitem[Ladegaard(2025)]{ladegaard2025}
Ladegaard, Isak. 2025. ``Differentiation by Disruption: Gatekeeper Perspectives on ‘AI-Aided Writing’ in Three Academic Disciplines.'' \emph{Socius}.

\bibitem[AI Shaming(2024)]{aishaming2024}
``AI Shaming: The Silent Stigma among Academic Writers and Researchers.'' \emph{Annals of Biomedical Engineering}, 2024.

\bibitem[Authenticity Paradox(2026)]{authenticity2026}
``The Authenticity Paradox: The Threshold Model of Synthetic Authenticity.'' \emph{International Journal of Human–Computer Interaction}, 2026.

\bibitem[Negotiating Authenticity(2024)]{negotiating2024}
``Negotiating the Authenticity of AI: How the Discourse on AI Rejects Human Indeterminacy.'' \emph{AI \& Society}, 2024.

\bibitem[Moral Hypocrites(2024)]{moral2024}
``Are Large Language Models Moral Hypocrites? A Study Based on Moral Foundations.'' arXiv:2405.11100, 2024.

\bibitem[Habermas Machine(2024)]{habermasmachine2024}
Google DeepMind. 2024. ``Habermas Machine'' for mediating online discourse.

\bibitem[Habermas Interview(2025)]{habermas2025}
Habermas, Jürgen. 2025. Interview with \emph{Süddeutsche Zeitung}, 16 April, objecting to the use of his name.

\bibitem[Mead(1934)]{mead1934}
Mead, George Herbert. 1934. \emph{Mind, Self, and Society}. Chicago: University of Chicago Press.

\bibitem[Parsons(1951)]{parsons1951}
Parsons, Talcott. 1951. \emph{The Social System}. Glencoe, IL: Free Press.

\bibitem[Adorno and Horkheimer(1947)]{adorno1947}
Adorno, Theodor W., and Max Horkheimer. 1947. \emph{Dialectic of Enlightenment}. Translated by John Cumming. New York: Continuum.

\bibitem[Marcuse(1964)]{marcuse1964}
Marcuse, Herbert. 1964. \emph{One-Dimensional Man}. Boston: Beacon Press.

\bibitem[Austin(1962)]{austin1962}
Austin, J. L. 1962. \emph{How to Do Things with Words}. Oxford: Clarendon Press.

\bibitem[Searle(1969)]{searle1969}
Searle, John R. 1969. \emph{Speech Acts: An Essay in the Philosophy of Language}. Cambridge: Cambridge University Press.

\bibitem[Lévi-Strauss(1950)]{levi1950}
Lévi-Strauss, Claude. 1950. ``Introduction à l'œuvre de Marcel Mauss.'' In \emph{Sociologie et Anthropologie}, by Marcel Mauss. Paris: Presses Universitaires de France.

\bibitem[Barthes(1957)]{barthes1957}
Barthes, Roland. 1957. \emph{Mythologies}. Paris: Seuil.

\bibitem[Chandler(2007)]{chandler2007}
Chandler, Daniel. 2007. \emph{Semiotics: The Basics}. 2nd ed. London: Routledge.

\bibitem[Stanley(2015)]{stanley2015}
Stanley, Jason. 2015. \emph{How Propaganda Works}. Princeton: Princeton University Press.

\bibitem[Cohen(1972)]{cohen1972}
Cohen, Stanley. 1972. \emph{Folk Devils and Moral Panics}. London: MacGibbon and Kee.

\bibitem[Goffman(1959)]{goffman1959}
Goffman, Erving. 1959. \emph{The Presentation of Self in Everyday Life}. Garden City, NY: Doubleday.

\bibitem[Bourdieu(1991)]{bourdieu1991}
Bourdieu, Pierre. 1991. \emph{Language and Symbolic Power}. Edited by John B. Thompson. Cambridge: Harvard University Press.

\bibitem[Foucault(1977)]{foucault1977}
Foucault, Michel. 1977. \emph{Discipline and Punish: The Birth of the Prison}. Translated by Alan Sheridan. New York: Pantheon.

\end{thebibliography}

\end{document}
